\section*{Source and counted claim}
Pierre-Antoine Guih\'eneuf, Patrice Le Calvez and Alejandro Passeggi, \emph{Area preserving homeomorphisms of surfaces with rational rotational direction}, Annales Henri Lebesgue 8 (2025), 329--372; exact official final native text processed 16 April 2025, \url{https://doi.org/10.5802/ahl.237}. Article mathematics and native diagrams are CC BY 4.0, \url{https://creativecommons.org/licenses/by/4.0/}. arXiv:2305.05755 is an identity alias only.

Only whole original Proposition 3.3 (\texttt{prop: nontransitivity3}) is counted. It is the authors' independent strengthened accumulation result, not the later area-preserving periodic-orbit theorem. The surface is connected, oriented, boundaryless and not the torus, with the non-singular oriented foliation and transverse-path conventions of the source. All three conclusions, including the accumulated boundary leaf and its side condition, are retained.

\paragraph{Complete current core.} The complete proof from the first paragraph through the final annulus and recurrence argument is preserved, including all nested lemmas, the full sublemma proof and all native TikZ figures. No support lemma is counted separately. The forward application pointer to a later dynamical strengthening is omitted as nondependency; the author's statement of Theorem 1.1 is retained only as evidence for another non-load-bearing forward reference.

\paragraph{Prerequisites and difficulty.} Jordan separation and ordinary surface covering theory are standard inputs. The source's elementary strip terminology and four endpoint alternatives remain as written. The H2 core constructs a transverse loop from recurrence, controls competing deck translates through local leaf-order and return configurations, excludes access from the wrong side, and proves that the projected saturated strip is an annulus. It does not assume the later forcing or periodic-orbit theorem. No proof is generated or repaired and no mathematical correctness certification is supplied.

\paragraph{Identity boundary.} Saved main and in-flight records have no exact DOI/arXiv/title hit. The geometric accumulation claim differs from mere existence of periodic points or homotopical rotation intervals. In-flight changes after 14 October 2026, 17:36:45 UTC remain unconfirmed; this is preparation only.
